\documentclass{article}
\usepackage[utf8]{inputenc}
\usepackage{tikz}
\title{Defi Lovelace N°1}
\author{Hugo Demaret}
\date{Octobre 2021}
\begin{document}
\maketitle
    \section*{Jeu en bois -}
        \textbf{Alice aime beaucoup jouer à un jeu qui contient des billes en bois colorés de deux couleurs différentes.}\\
        \textbf{Ce jeu se joue sur une planche en bois avec n trous, qui sont tous remplis d'une seule bille (il y a donc n billes). \\Une bille peut être de deux couleurs : ici bleue ou rouge.}\\\\
        \textit{Voici un schéma du jeu en question avec 8 billes :}
        \begin{center}
            \begin{table}[h]
                \begin{tabular}{|c|c|c|c|c|c|c|c| p{15cm}}
                    \hline 
                    \begin{tikzpicture}
                        \node[circle,draw=red,fill=red,minimum size = 25](a) at (0,0) {};
                    \end{tikzpicture}&
                    \begin{tikzpicture}
                        \node[circle,draw=blue,fill=blue,minimum size = 25](b) at (1,0) {};
                    \end{tikzpicture}&
                    \begin{tikzpicture}
                        \node[circle,draw=blue,fill=blue,minimum size = 25](c) at (2,0) {};
                    \end{tikzpicture}&
                    \begin{tikzpicture}
                        \node[circle,draw=blue,fill=blue,minimum size = 25](d) at (3,0) {};
                    \end{tikzpicture}&
                    \begin{tikzpicture}
                        \node[circle,draw=red,fill=red,minimum size = 25](e) at (4,0) {};
                    \end{tikzpicture}&
                    \begin{tikzpicture}
                        \node[circle,draw=red,fill=red,minimum size = 25](f) at (5,0) {};
                    \end{tikzpicture}&
                    \begin{tikzpicture}
                        \node[circle,draw=blue,fill=blue,minimum size = 25](g) at (6,0) {};
                    \end{tikzpicture}&
                    \begin{tikzpicture}
                        \node[circle,draw=red,fill=red,minimum size = 25](h) at (7,0) {};
                    \end{tikzpicture}\\
                    \hline
                \end{tabular}
            \end{table}
        \end{center}
        \textbf{Le but du jeu est de séparer les billes bleues des billes rouges, de sorte à ce qu'un coté (arbitrairement choisit) contienne les billes rouges, et l'autre les billes bleues.}
        \subsection*{1 -}
            \textsf{Proposez une solution à ce problème, sous la forme d'un script python.}\\
        \subsection*{2 -}
            \textsf{Ryan, un ami d'Alice, dit que pour un jeu de n billes, on peut toujours résoudre la planche en au plus n+1 coups. A-t-il raison ?}\\
            \textit{Expliquez votre raisonnement.}
        \subsection*{Participation :}
            \textsf{Vous pouvez réaliser ce défi quand vous voulez.}\\
            \textsf{Le premier à le résoudre sera le lauréat, et sa copie sera (avec son accord) publiée avec la correction.}\\
            \textit{Pour le rendre, envoyez un mail à hugo.demaret(at)defi-lovelace.fr avec pour objet [DEFI 1]. \\ Mettez dans votre mail votre solution le .py de votre code python, ainsi que les réponses aux questions dans un pdf.}\\
            \textit{Précisez votre nom (si vous le souhaitez) afin que l'on puisse revenir vers vous.}
            \textbf{Bon courage :)}
\end{document}